\songtitle{Пѣснь о Кругахъ Путей}

\tag*{2026}\tag{roads}
\begin{notes}
Lyrics by Grok based in the introductory summary by Claude of the Circuits and Roads document.
\end{notes}
\begin{style}
nu metal, industrial rock, post-punk blues, ancient epic bridges, throat singing chorus, mezzo vocal lead, Ruthenian language vocals, gang shouts, bowed low strings, distorted baritone guitar, detuned six-string riff, stomp snare march, tom-heavy breakdowns, church bell drones, tape saturation, analog room mics, plate reverb, dark heroic, 112 BPM, ritual march
\end{style}
\begin{lyrics}
\songpart{Verse 1}
Въ лѣтѣ тысяча девятьсотъ пятнадцать,
Списки мѣстъ, якъ сны нерѣшенны,
Города и храмы, древни слои,
Стамбулъ, Петра, Алгамбра, Анґкоръ,
Таджъ Махалъ, Лима, Мѣксико —
Магниты душѣ, любопытства зовъ.

\songpart{Chorus}
Не просто мѣста — а круги пути,
Не „куда поѣхать“, а „системы знати“.
Географiя, исторiя, рѣки и торгъ,
Тридцать днiй въ каждомъ — жизнь проектъ великъ.
Десять круговъ, триста днiй странствiя,
Черезъ годы, черезъ судьбы — завершити.

\songpart{Verse 2}
Цивилизацiи слои накоплены,
Великiи грады — Лондонъ, Прага, Сеулъ,
Чудеса природы — Патагонiя, Виктория,
Желтоустонъ, Кано Кристалесъ зовутъ.
Списки стары, почти безъ перемѣны,
Но духъ измѣнился — отъ желанiй къ пониманiю.

\songpart{Bridge}
Географiческа логика — что земля даетъ,
Историческiй разсказъ — слои вѣковъ,
Тридцатидневный путь съ глубиной временъ,
И гастрономiя — вкусъ земли и предковъ.
Балтика–Валдай, Шелковый Путь —
Геополитика ждетъ, но мечта живетъ.

\songpart{Final Chorus (with variation)}
Изъ списковъ 2015 — въ карты любопытства,
Въ круги системъ, въ пониманiе свѣта.
Не праздникъ короткiй — проектъ жизни всей,
Въ каждомъ мѣстѣ — связь съ прошлымъ и съ собой.
Русь моя душа, путь вѣчный зоветъ,
Кругами свѣта — до конца вѣковъ!
\end{lyrics}

